\subsection{Semilinear involutions}

Let E be a complex vector space. A semilinear involution on E is an R-linear map from E to E such that \(u \circ u = Id_{E}\) and \(u(\lambda x) = \overline{\lambda}u(x)\) for all \(\lambda \in C\) and every \(x \in E\) . We then denote by \(E^{u}\) the real vector subspace of E formed by the elements \(x \in E\) such that \(u(x) = x\) .

Lemma 1. --- Let E be a complex vector space and let u be a semilinear involution on E. Let \(x \in \mathrm{E}\) ; set

\[
x _ {1} = \frac {1}{2} (x + u (x)), \quad x _ {2} = \frac {1}{2 i} (x - u (x)).
\]

The pair \((x_{1}, x_{2})\) is the unique element of \(E^{u} \times E^{u}\) such that \(x = x_{1} + ix_{2}\) . The elements \(x_{1}\) and \(x_{2}\) satisfy \(x_{1} + ix_{2} = x\) and belong to \(E^{u}\) since \(u(u(x)) = x\) . Conversely, if \(y_{1}\) and \(y_{2}\) are in \(E^{u}\) and satisfy \(x = y_{1} + iy_{2}\) , it follows that \(u(x) = u(y_{1}) + u(iy_{2}) = y_{1} - iy_{2}\) , hence

\[
y _ {1} = \frac {1}{2} (x + u (x)) = x _ {1}, \quad i y _ {2} = \frac {1}{2} (x - u (x)) = i x _ {2}.
\]

PROPOSITION 1. --- Let \(E_{1}\) and \(E_{2}\) be complex vector spaces, and let \(u_{1}\) and \(u_{2}\) be semilinear involutions on \(E_{1}\) and \(E_{2}\) respectively. Let f be a linear map from \(E_{1}\) into \(E_{2}\) . Then \(f \circ u_{1} = u_{2} \circ f\) if and only if \(f(\mathrm{E}_{1}^{u_{1}}) \subset \mathrm{E}_{2}^{u_{2}}\) .

If \(f \circ u_1 = u_2 \circ f\) , we immediately obtain \(f(\mathrm{E}_1^{u_1}) \subset \mathrm{E}_2^{u_2}\) . Conversely, suppose this condition is satisfied. Let \(x \in \mathrm{E}\) . Write \(x = x_1 + ix_2\) with \((x_1, x_2) \in \mathrm{E}_1^{u_1} \times \mathrm{E}_1^{u_1}\) (lemma above). We then have \(f(u_1(x)) = f(x_1) - if(x_2)\) and \(u_2(f(x)) = u_2(f(x_1)) - iu_2(f(x_2)) = f(x_1) - if(x_2) = f(u_1(x))\) .

